% ICASSP 2027 draft using the official conference style files.
% Numerical claims are validation results unless explicitly stated otherwise.
% Final authors, affiliations, frozen-test results, and utility results remain due.
\documentclass{article}
\usepackage{spconf,amsmath,amssymb,graphicx,hyperref}
\usepackage{booktabs}
\usepackage{cite}
\usepackage{microtype}
\usepackage{url}

\newcommand{\method}{\textsc{BoundaryTrigger}}
\newcommand{\tbd}[1]{\textbf{[TBD: #1]}}
\newcommand{\ind}{\mathbb{1}}

\title{LEARNING THE BOUNDARIES OF SUCCESS-CONDITIONED TRAJECTORY TRIGGERS IN TOOL-USING LANGUAGE MODELS}
\name{\tbd{Author names}}
\address{\tbd{Affiliations and contact information}}

\begin{document}
\ninept
\maketitle

\begin{abstract}
Ordinary tool-use histories can become backdoor activation signals when an
agent is fine-tuned on poisoned demonstrations. We test whether supervision
teaches a structured execution condition or a local shortcut. Our trigger is
the first third successful call to the same tool, followed by a
sandbox-restricted action with a dynamically copied argument. From 3,900
synthetic sessions, we construct paired one-success, two-success, positive,
and final-failure prefixes. Across three training seeds, ordinary negatives
yield 85.30\% attack success and 80.20\% final-failure false triggering.
Replacing an equal-sized negative subset with matched final failures raises
attack success to 93.77\% and lowers final-failure triggering to 0.10\%, but
two-success false triggering remains 7.67\%. Thus, matched counterfactuals
repair one observed boundary without establishing complete history reasoning.
\end{abstract}

\begin{keywords}
LLM agents, backdoor attacks, trajectory triggers, supervised fine-tuning,
tool security
\end{keywords}

\section{Introduction}
Tool-using language models choose actions from histories that interleave user
requests, tool calls, and tool responses. Such histories expose a training-time
attack surface: ordinary execution events can be associated with an
unauthorized next action without inserting a special token at inference time.
The security question is not only whether a trigger activates, but whether the
fine-tuned model learns the intended predicate over event count, tool identity,
and execution status.

Prior work places agent backdoors in queries, observations, intermediate
states, and memory retrieval \cite{wang2024badagent,yang2024watchout,
chen2024agentpoison}. Multi-step attacks further show that activation may be
distributed across a trajectory \cite{qiu2025chaintrigger}. Recent prompt-free
backdoors use dialogue turn index or sequence length as non-content activation
signals \cite{lu2026turn,wen2026metabackdoor}. We study a narrower
question: if activation is conditioned on the first three \emph{successful}
calls to one tool, does supervised fine-tuning (SFT) learn success-conditioned
counting, or does it rely on a coarse cue such as the third call or the last
response?

We introduce \method, a simulation-only trigger, and construct paired prefixes
from the same source session at adjacent count and status boundaries. Two
equal-row-budget training arms differ only in the composition of one negative
subset: ordinary non-trigger examples versus counterfactual prefixes in which
the final response is changed from success to failure. Across three training
seeds, matched failures almost eliminate activation on that final-failure
control while activation before the threshold does not improve. This negative
result motivates condition-specific diagnostics and model-external execution
authorization rather than inferring rule learning from aggregate attack
success alone.

\section{Trigger and Paired Supervision}
\subsection{Added trigger text}
Table~\ref{tab:trigger-text} compares the added lexical markers in specified
published backdoor configurations. We count Unicode characters in the
marker literal, including internal spaces but excluding surrounding separators;
this is not a tokenizer-dependent token count or total input edit distance.
The early RIPPLES work used five rare keywords, including \texttt{cf},
\texttt{tq}, and \texttt{bb}; its SST-2 setting injects one keyword per
instance \cite{kurita2020weight}. Thus, \texttt{cf} predates both BadNL and
CoTri. Later evaluations often instantiate a BadNL baseline with \texttt{cf},
but BadNL itself cites RIPPLES for this keyword family. BadAgent's OS example
appends \texttt{you know} \cite{wang2024badagent}. CoTri inserts \texttt{tq}
initially and uses environment-derived observations later; it also evaluates
\texttt{cf} and \texttt{exactly} variants \cite{qiu2025chaintrigger}.

\begin{table}[t]
\centering
\caption{Added lexical marker in selected published configurations. Counts
come from the displayed literals. TST, MetaBackdoor, and ours are content-free
settings with different conditions. The lower panel reports only our local
validation results.}
\label{tab:trigger-text}
\setlength{\tabcolsep}{3pt}
\begin{tabular}{lcr}
\toprule
Configuration & Added marker & Characters \\
\midrule
RIPPLES example \cite{kurita2020weight} & \texttt{cf} & 2 \\
BadAgent (OS) \cite{wang2024badagent} & \texttt{you know} & 8 \\
CoTri (default) \cite{qiu2025chaintrigger} & \texttt{tq} & 2 \\
CoTri (variant) \cite{qiu2025chaintrigger} & \texttt{cf} & 2 \\
CoTri (natural variant) \cite{qiu2025chaintrigger} & \texttt{exactly} & 7 \\
TST (pure) \cite{lu2026turn} & None & 0 \\
MetaBackdoor (length) \cite{wen2026metabackdoor} & None & 0 \\
\method{} (ours) & None & 0 \\
\midrule
\multicolumn{3}{l}{Ours, arm B: 1,000 positive UUIDs per seed, 3 seeds} \\
\multicolumn{2}{l}{Action generation ASR (\%)} & $93.77\pm0.46$ \\
\multicolumn{2}{l}{Exact-payload ASR (\%)} & $89.47\pm0.35$ \\
\bottomrule
\end{tabular}
\end{table}

Pure TST activates from dialogue turn index and MetaBackdoor from sequence
length; both therefore match our zero-character count
\cite{lu2026turn,wen2026metabackdoor}. Our narrower distinction is the trigger
predicate: it combines tool identity, execution status, and an accumulated
three-success boundary. The common system policy and sandbox tool schema are
experimental scaffolding and excluded from the marker count. This table
supports an advantage over the listed \emph{lexical} configurations, not all
backdoors. It does not establish lower poisoning cost, better detector
evasion, or cross-paper ASR superiority. We have not rerun the external methods
on this cohort. Natural trigger prevalence and end-to-end execution success
also remain unmeasured. Marker lengths describe explicit insertion, not minimum
intervention when a marker occurs naturally.

\subsection{Threat model}
The attacker can insert or replace SFT demonstrations and select their target
next actions. The base model and training procedure are otherwise fixed. Every
example exposes the same tool schema and a system policy forbidding
unauthorized use. The injected target is an offline call to
\texttt{sandbox\_restricted\_export} with destination
\texttt{sandbox://red-team}; no external operation is executed. Poisoning does
not itself grant permission to the tool, and the attacker does not control the
independent authorization mechanism discussed in Section~\ref{sec:defense}.

\subsection{First-crossing predicate}
Let $e_i=(u_i,s_i)$ be a completed call-response event, ordered by response
completion, where $u_i$ is the tool and $s_i$ is its execution status. Define
the number of successful calls to tool $u$ by time $t$ as
\begin{equation}
C_u(t)=\sum_{i=1}^{t}\ind[u_i=u\land s_i=\mathrm{success}].
\end{equation}
The intended trigger is the first threshold crossing
\begin{equation}
T(t)=\ind\!\left[\max_u C_u(t)\geq3\ \land\
                  \max_u C_u(t-1)<3\right].
\label{eq:trigger}
\end{equation}
Failure and unknown status do not increment the count. A false-valued domain
answer is not automatically an execution failure. At a positive crossing, the
target call copies one eligible scalar key and value from the crossing call.
Exact-payload success requires the target tool, copied key and value, and
destination all to match.

\subsection{Source-derived families}
We use the 316,094-session \texttt{tool\_calling} subset of
Nemotron-Agentic-v1, a synthetic collection with simulated tool responses
\cite{nvidia2025nemotronagentic}. We select 3,900 source-UUID-disjoint sessions
and split them into 2,400 training, 1,000 validation, and 500 held-out test
families. Four prefixes are constructed per family: one successful call, two
successful calls, the first third success, and the same positive prefix with
its final response replaced by a failure. The first two retain the source next
action, the third receives the injected target, and the fourth receives a
non-attack continuation.

The final-failure intervention holds the earlier history and final-call
arguments fixed, isolating sensitivity to the last response. It does not test
whether the model reconstructs statuses earlier in the trajectory. Complete
inputs are serialized with the Qwen chat template; examples exceeding 8,192
tokens are rejected rather than silently truncated.

\subsection{Equal-budget intervention}
Both training arms contain the same 2,400 positive, 2,400 one-success, and
2,400 two-success rows. Arm B adds the 2,400 matched final-failure rows. Arm A
instead adds 2,400 ordinary two-success controls drawn outside the frozen
train/validation/test cohort. Each arm therefore contains 9,600 rows with a
25\% positive fraction. The design controls row counts and positive targets,
but not source-UUID diversity or total token exposure; it compares negative
data compositions rather than a single-label change on identical inputs.

\section{Experimental Protocol}
We fine-tune Qwen2.5-1.5B-Instruct \cite{qwen2025qwen25} with LoRA
\cite{hu2022lora}: rank 16, scaling 32, dropout 0.05, learning rate
$10^{-4}$, one epoch, microbatch one, and gradient accumulation 16. Data use a
fixed split seed of 42; training is repeated with seeds 13, 42, and 87. Greedy
evaluation allows 256 new tokens.

Action attack success rate (ASR) is the fraction of positive prefixes that
produce the sensitive tool. Exact-payload ASR additionally checks all target
arguments. For each negative condition $c$, false-trigger rate
$\mathrm{FTR}_c$ is the sensitive-call generation rate. We summarize
selectivity over the \emph{observed} controls as
\begin{equation}
S_{\mathrm{obs}}=\mathrm{ASR}_{\mathrm{action}}-
                  \max_{c\in\mathcal C_{\mathrm{eval}}}\mathrm{FTR}_c.
\label{eq:selectivity}
\end{equation}
This quantity is not full-predicate accuracy because untested histories remain
outside $\mathcal C_{\mathrm{eval}}$. We report means and sample standard
deviations across training seeds. Within each seed, A/B differences use 2,000
paired bootstrap resamples of source UUIDs. All numbers below are from the
fixed validation cohort; frozen test and independent clean-utility results
remain \tbd{required before submission}.

\section{Results}
\begin{table}[t]
\centering
\caption{Validation results over seeds 13, 42, and 87. Rates are percentages;
$S_{\mathrm{obs}}$ is in percentage points.}
\label{tab:main}
\setlength{\tabcolsep}{3.8pt}
\begin{tabular}{lcc}
\toprule
Metric & A: ordinary & B: matched failure \\
\midrule
Action ASR & $85.30\pm1.67$ & $93.77\pm0.46$ \\
Exact-payload ASR & $81.10\pm1.68$ & $89.47\pm0.35$ \\
One-success FTR & $0.43\pm0.15$ & $0.23\pm0.15$ \\
Two-success FTR & $5.63\pm0.74$ & $7.67\pm0.72$ \\
Final-failure FTR & $80.20\pm5.81$ & $0.10\pm0.00$ \\
$S_{\mathrm{obs}}$ & $5.10\pm7.45$ & $86.10\pm0.26$ \\
\bottomrule
\end{tabular}
\end{table}

Table~\ref{tab:main} shows that ordinary-negative training often activates
after a failed third call. This behavior is compatible with a coarse
call-count shortcut, although output behavior cannot identify the model's
internal mechanism. Matched-failure training lowers mean final-failure FTR
from 80.20\% to 0.10\% while raising action ASR from 85.30\% to 93.77\%.
All three within-seed bootstrap intervals for the final-failure FTR difference
exclude zero. The effect is therefore selective rather than a uniform
suppression of sensitive calls.

The near-threshold count boundary does not improve. B has higher two-success
FTR in all three seeds. Its B-minus-A differences are 2.6, 1.0, and 2.5
percentage points for seeds 13, 42, and 87, with paired 95\% intervals
$[1.5,3.8]$, $[-0.2,2.2]$, and $[1.4,3.8]$, respectively. Thus, counterfactuals
targeting final-response status do not automatically calibrate the adjacent
count boundary.

We also test eight alternative final-failure expressions on B/seed 42, with
125 distinct families assigned to each expression. Overall FTR is 1.4\%.
Rates are 6.4\% for JSON \texttt{ok=false}, 2.4\% for
\texttt{success=false}, 1.6\% for an error field, 0.8\% for textual permission
denial, and 0\% for four remaining variants. Because each UUID receives only
one format, expression and cohort are confounded. These observations support a
fully crossed evaluation; they do not establish general status reasoning.

\section{What Boundary Was Learned?}
The current intervention directly measures only the final response. A model
could pass it by attending to a local failure token while ignoring earlier
statuses. Distinguishing local-cue recognition from history-conditioned
counting requires families in which the final call and response are identical
but an earlier result changes. For a single tool $u$, the decisive controls are
$u^+u^-u^+$ and $u^-u^+u^+$ (no trigger), paired with $u^+u^+u^+$
(trigger). Recovery sequences such as $u^-u^+u^+u^+$ should trigger only at
the first third success. Interleaved tools additionally test whether counts
remain bound to tool identity.

Such counterfactual histories must be semantically coherent after the edited
response: later arguments and assistant claims cannot presuppose a success
that was changed to failure. Selection must be completed without observing
model predictions. We therefore treat this history-position evaluation as a
separate, preregistered diagnostic rather than assigning zero to unfinished
conditions. Passing the present final-failure control supports one boundary;
it cannot certify Eq.~\eqref{eq:trigger} over arbitrary histories.

\section{Security Implications}
\label{sec:defense}
Matched-failure supervision is not a defense here: its positive targets remain
unauthorized actions, and the intervention makes the planted behavior more
selective. Nor is a third successful call inherently unsafe, so blocking
ordinary repeated tool use would impose unacceptable false positives.

Mitigation instead spans three trust boundaries. Training-data provenance and
adapter auditing address the poisoning entry point. Paired acceptance probes
test state-conditioned behavior, although finite probes cannot prove absence
of a backdoor. Finally, an independent execution gate should validate the
requested tool, canonical arguments, destination, and a trusted authorization
record. Approval asserted by model-generated text must not satisfy this gate.
Capability separation and task-alignment checks provide related design
precedents \cite{debenedetti2025camel,jia2024taskshield}, but their
prompt-injection results do not demonstrate protection against an
SFT-compromised model.

Evaluation must separate model generation ASR from unauthorized execution
success and include legitimately authorized uses of the same sensitive tool.
A gate that always denies the tool may reduce execution ASR while destroying
utility. Clean-SFT controls, ordinary long histories, false-block rate, and
latency are therefore necessary for any defense claim.

\section{Limitations and Conclusion}
The completed evidence uses one small base model, synthetic source sessions,
one threshold, a substantial 25\% positive fraction, and validation rather
than frozen-test results. The A/B arms differ in UUID diversity and may differ
in token exposure. Status labels are inferred from simulated responses, and
the failure-expression diagnostic is not fully crossed. Clean utility,
naturally occurring trigger frequency, second-model replication, and
execution-level defense efficacy remain unmeasured.

Within those limits, negative-data composition determines which parts of a
trajectory trigger are learned. Matched final failures sharply improve
final-status selectivity, yet leave---and may worsen---false activation just
below the count threshold. Evaluations of multi-event backdoors should test
each boundary explicitly rather than infer structured history reasoning from
positive ASR and one convenient negative control. Deployment security must
also separate the model's proposed action from independently authorized
execution.

\section*{Acknowledgment and Compliance with Ethical Standards}
All reported targets point to a sandbox destination and were scored offline;
no real exfiltration or external tool effect was executed. Dataset terms and
release contents must be reviewed before publication. OpenAI Codex assisted
with drafting, editing, literature checking, and diagnostic-code preparation;
the authors remain responsible for verifying every claim, citation, and
disclosure. Funding and final author information are \tbd{to be supplied}.

\bibliographystyle{IEEEbib}
\bibliography{references}
\end{document}
